\section{Data}
\label{sec:data}
The UCI file \cite{lohweg2012uci} (CC BY 4.0) has 1,372 rows. Each has four inputs, the variance, skewness, kurtosis and entropy of a wavelet-transformed image, and a binary class. There are no missing values. We remove 24 exact duplicate rows, all from class 0, which leaves 1,348 notes, 45.3\% of them in class 1 (Table~\ref{tab:data}). UCI does not say which class is forged, and secondary sources disagree, so we call class 1 ``forged''. Error and AUC do not depend on this choice; missed forgeries and false alarms would swap.

\begin{table}[t]
  \centering
  \caption{Data size and shape at each stage. ``Forged'' is the share of class 1.}
  \label{tab:data}
  \small
  % generated by src/make_tables.py
\begin{tabular}{lrrr}
\toprule
Stage & Rows & Inputs & Forged \\
\midrule
Raw UCI file & 1372 & 4 & 44.5\% \\
After removing exact duplicates & 1348 & 4 & 45.3\% \\
Development set (80\%) & 1078 & 4 & 45.3\% \\
Test set (20\%, held out) & 270 & 4 & 45.2\% \\
\bottomrule
\end{tabular}
% 24 duplicate rows dropped

\end{table}

We make one stratified 80/20 split with a fixed random seed\footnote{The seed is 5420, the course number. It is only a fixed starting point for the random-number generator, so that anyone rerunning the code gets the same split; it was not chosen by looking at results. Every later random choice derives its seed from it.} into 1,078 development notes and 270 test notes, and store it with the data. All tuning, learning curves and model choices use the development notes only. The test notes were used once, after every modelling choice was fixed (Section~\ref{sec:test}).

\paragraph{Several images per specimen?} The file has no specimen identifiers, but the class-1 rows are strongly correlated in file order: the lag-1 autocorrelation of skewness is 0.62 for class 1 and about 0 for class 0, with a repeating pattern roughly every 8 rows. This suggests that the forged class contains several images of each physical specimen in a row. If near-identical images of one note fall on both sides of a split, a method that memorises training points, such as 1-nearest-neighbour, gains an unfair advantage. We cannot remove this risk without identifiers, so we test for it: we treat each run of 8 consecutive rows of one class as a pseudo-specimen and repeat the key experiments with whole pseudo-specimens kept on one side of every split (Section~\ref{sec:grouped}).

\section{Method}
\label{sec:method}
\paragraph{Classifiers.} All inputs are standardised inside each training set (a scikit-learn pipeline \cite{pedregosa2011scikit}), so no statistics leak from evaluation rows. Each method has a small grid tuned by stratified 5-fold CV on the training set alone:
\begin{itemize}
  \item \textbf{Logistic regression} (baseline): $C\in\{10^{-3},10^{-2.5},\dots,10^{3}\}$, 13 values.
  \item \textbf{$k$NN}: $k\in\{1,3,5,7,9,11,15,21,31\}$, limited to $k$ smaller than an inner training fold.
  \item \textbf{Random forest}$^\dagger$: 200 trees; features per split $m\in\{1,2,4\}$, minimum leaf size $\in\{1,3\}$.
  \item \textbf{Gradient boosting}$^\dagger$: 200 trees; learning rate $\eta\in\{0.05,0.2\}$, depth $\in\{1,3\}$.
  \item \textbf{Neural network}$^\dagger$: one hidden layer with 8 or 32 ReLU units, weight decay $\alpha\in\{10^{-4},10^{-2},1\}$, trained by L-BFGS.
\end{itemize}
$^\dagger$Later-course extensions (ensembles and neural networks). No baseline number depends on them. The grids are deliberately small so that every method can be tuned inside each of 2,750 training draws. Appendix~\ref{app:forest} shows that a 64-setting forest grid on all development data does not do better than the small one.

\paragraph{Baseline estimate.} On all 1,078 development notes we report the 5-fold CV error of logistic regression at the chosen $C$, with the standard error (SE) of the five fold errors, SD$/\sqrt{5}$. Because $C$ is chosen on the same folds, we also report nested 5$\times$5 CV.

\paragraph{Learning curves.} For each training size $n\in\{10,20,30,50,100,150,200,300,400,600,800\}$ and each of 50 draws, we sample $n$ development notes stratified by class, tune each method by 5-fold CV inside the sample, and score it on all remaining development notes (between 278 and 1,068). All methods see the same 50 training sets at each size, so we compare them with paired differences: for method $j$ and draw $d$, $\Delta_{jd}=\text{err}_{jd}-\text{err}_{\text{LR},d}$, reported as the mean over draws $\pm$ its SE. We call a difference clear when it exceeds 2 SE.

\paragraph{Notes needed to reach a target.} For each method we find where the mean curve first falls to 2\% and to 1\% error, by linear interpolation between neighbouring sizes on log--log axes. We get a 95\% interval by resampling the 50 draws at every size with replacement (2,000 bootstrap replicates) and repeating the calculation.

\paragraph{Test evaluation.} Each method is tuned by 5-fold CV on all development notes, refitted on them, and scored once on the 270 test notes. We give Wilson 95\% intervals for each error, and compare each method with the baseline by an exact McNemar test on the notes where the two disagree \cite{dietterich1998approximate}.
